% Options for packages loaded elsewhere
% Options for packages loaded elsewhere
\PassOptionsToPackage{unicode}{hyperref}
\PassOptionsToPackage{hyphens}{url}
\PassOptionsToPackage{dvipsnames,svgnames,x11names}{xcolor}
%
\documentclass[
  letterpaper,
  DIV=11,
  numbers=noendperiod]{scrartcl}
\usepackage{xcolor}
\usepackage{amsmath,amssymb}
\setcounter{secnumdepth}{-\maxdimen} % remove section numbering
\usepackage{iftex}
\ifPDFTeX
  \usepackage[T1]{fontenc}
  \usepackage[utf8]{inputenc}
  \usepackage{textcomp} % provide euro and other symbols
\else % if luatex or xetex
  \usepackage{unicode-math} % this also loads fontspec
  \defaultfontfeatures{Scale=MatchLowercase}
  \defaultfontfeatures[\rmfamily]{Ligatures=TeX,Scale=1}
\fi
\usepackage{lmodern}
\ifPDFTeX\else
  % xetex/luatex font selection
\fi
% Use upquote if available, for straight quotes in verbatim environments
\IfFileExists{upquote.sty}{\usepackage{upquote}}{}
\IfFileExists{microtype.sty}{% use microtype if available
  \usepackage[]{microtype}
  \UseMicrotypeSet[protrusion]{basicmath} % disable protrusion for tt fonts
}{}
\makeatletter
\@ifundefined{KOMAClassName}{% if non-KOMA class
  \IfFileExists{parskip.sty}{%
    \usepackage{parskip}
  }{% else
    \setlength{\parindent}{0pt}
    \setlength{\parskip}{6pt plus 2pt minus 1pt}}
}{% if KOMA class
  \KOMAoptions{parskip=half}}
\makeatother
% Make \paragraph and \subparagraph free-standing
\makeatletter
\ifx\paragraph\undefined\else
  \let\oldparagraph\paragraph
  \renewcommand{\paragraph}{
    \@ifstar
      \xxxParagraphStar
      \xxxParagraphNoStar
  }
  \newcommand{\xxxParagraphStar}[1]{\oldparagraph*{#1}\mbox{}}
  \newcommand{\xxxParagraphNoStar}[1]{\oldparagraph{#1}\mbox{}}
\fi
\ifx\subparagraph\undefined\else
  \let\oldsubparagraph\subparagraph
  \renewcommand{\subparagraph}{
    \@ifstar
      \xxxSubParagraphStar
      \xxxSubParagraphNoStar
  }
  \newcommand{\xxxSubParagraphStar}[1]{\oldsubparagraph*{#1}\mbox{}}
  \newcommand{\xxxSubParagraphNoStar}[1]{\oldsubparagraph{#1}\mbox{}}
\fi
\makeatother

\usepackage{color}
\usepackage{fancyvrb}
\newcommand{\VerbBar}{|}
\newcommand{\VERB}{\Verb[commandchars=\\\{\}]}
\DefineVerbatimEnvironment{Highlighting}{Verbatim}{commandchars=\\\{\}}
% Add ',fontsize=\small' for more characters per line
\usepackage{framed}
\definecolor{shadecolor}{RGB}{241,243,245}
\newenvironment{Shaded}{\begin{snugshade}}{\end{snugshade}}
\newcommand{\AlertTok}[1]{\textcolor[rgb]{0.68,0.00,0.00}{#1}}
\newcommand{\AnnotationTok}[1]{\textcolor[rgb]{0.37,0.37,0.37}{#1}}
\newcommand{\AttributeTok}[1]{\textcolor[rgb]{0.40,0.45,0.13}{#1}}
\newcommand{\BaseNTok}[1]{\textcolor[rgb]{0.68,0.00,0.00}{#1}}
\newcommand{\BuiltInTok}[1]{\textcolor[rgb]{0.00,0.23,0.31}{#1}}
\newcommand{\CharTok}[1]{\textcolor[rgb]{0.13,0.47,0.30}{#1}}
\newcommand{\CommentTok}[1]{\textcolor[rgb]{0.37,0.37,0.37}{#1}}
\newcommand{\CommentVarTok}[1]{\textcolor[rgb]{0.37,0.37,0.37}{\textit{#1}}}
\newcommand{\ConstantTok}[1]{\textcolor[rgb]{0.56,0.35,0.01}{#1}}
\newcommand{\ControlFlowTok}[1]{\textcolor[rgb]{0.00,0.23,0.31}{\textbf{#1}}}
\newcommand{\DataTypeTok}[1]{\textcolor[rgb]{0.68,0.00,0.00}{#1}}
\newcommand{\DecValTok}[1]{\textcolor[rgb]{0.68,0.00,0.00}{#1}}
\newcommand{\DocumentationTok}[1]{\textcolor[rgb]{0.37,0.37,0.37}{\textit{#1}}}
\newcommand{\ErrorTok}[1]{\textcolor[rgb]{0.68,0.00,0.00}{#1}}
\newcommand{\ExtensionTok}[1]{\textcolor[rgb]{0.00,0.23,0.31}{#1}}
\newcommand{\FloatTok}[1]{\textcolor[rgb]{0.68,0.00,0.00}{#1}}
\newcommand{\FunctionTok}[1]{\textcolor[rgb]{0.28,0.35,0.67}{#1}}
\newcommand{\ImportTok}[1]{\textcolor[rgb]{0.00,0.46,0.62}{#1}}
\newcommand{\InformationTok}[1]{\textcolor[rgb]{0.37,0.37,0.37}{#1}}
\newcommand{\KeywordTok}[1]{\textcolor[rgb]{0.00,0.23,0.31}{\textbf{#1}}}
\newcommand{\NormalTok}[1]{\textcolor[rgb]{0.00,0.23,0.31}{#1}}
\newcommand{\OperatorTok}[1]{\textcolor[rgb]{0.37,0.37,0.37}{#1}}
\newcommand{\OtherTok}[1]{\textcolor[rgb]{0.00,0.23,0.31}{#1}}
\newcommand{\PreprocessorTok}[1]{\textcolor[rgb]{0.68,0.00,0.00}{#1}}
\newcommand{\RegionMarkerTok}[1]{\textcolor[rgb]{0.00,0.23,0.31}{#1}}
\newcommand{\SpecialCharTok}[1]{\textcolor[rgb]{0.37,0.37,0.37}{#1}}
\newcommand{\SpecialStringTok}[1]{\textcolor[rgb]{0.13,0.47,0.30}{#1}}
\newcommand{\StringTok}[1]{\textcolor[rgb]{0.13,0.47,0.30}{#1}}
\newcommand{\VariableTok}[1]{\textcolor[rgb]{0.07,0.07,0.07}{#1}}
\newcommand{\VerbatimStringTok}[1]{\textcolor[rgb]{0.13,0.47,0.30}{#1}}
\newcommand{\WarningTok}[1]{\textcolor[rgb]{0.37,0.37,0.37}{\textit{#1}}}

\usepackage{longtable,booktabs,array}
\usepackage{calc} % for calculating minipage widths
% Correct order of tables after \paragraph or \subparagraph
\usepackage{etoolbox}
\makeatletter
\patchcmd\longtable{\par}{\if@noskipsec\mbox{}\fi\par}{}{}
\makeatother
% Allow footnotes in longtable head/foot
\IfFileExists{footnotehyper.sty}{\usepackage{footnotehyper}}{\usepackage{footnote}}
\makesavenoteenv{longtable}
\usepackage{graphicx}
\makeatletter
\newsavebox\pandoc@box
\newcommand*\pandocbounded[1]{% scales image to fit in text height/width
  \sbox\pandoc@box{#1}%
  \Gscale@div\@tempa{\textheight}{\dimexpr\ht\pandoc@box+\dp\pandoc@box\relax}%
  \Gscale@div\@tempb{\linewidth}{\wd\pandoc@box}%
  \ifdim\@tempb\p@<\@tempa\p@\let\@tempa\@tempb\fi% select the smaller of both
  \ifdim\@tempa\p@<\p@\scalebox{\@tempa}{\usebox\pandoc@box}%
  \else\usebox{\pandoc@box}%
  \fi%
}
% Set default figure placement to htbp
\def\fps@figure{htbp}
\makeatother





\setlength{\emergencystretch}{3em} % prevent overfull lines

\providecommand{\tightlist}{%
  \setlength{\itemsep}{0pt}\setlength{\parskip}{0pt}}



 


\KOMAoption{captions}{tableheading}
\makeatletter
\@ifpackageloaded{caption}{}{\usepackage{caption}}
\AtBeginDocument{%
\ifdefined\contentsname
  \renewcommand*\contentsname{Table of contents}
\else
  \newcommand\contentsname{Table of contents}
\fi
\ifdefined\listfigurename
  \renewcommand*\listfigurename{List of Figures}
\else
  \newcommand\listfigurename{List of Figures}
\fi
\ifdefined\listtablename
  \renewcommand*\listtablename{List of Tables}
\else
  \newcommand\listtablename{List of Tables}
\fi
\ifdefined\figurename
  \renewcommand*\figurename{Figure}
\else
  \newcommand\figurename{Figure}
\fi
\ifdefined\tablename
  \renewcommand*\tablename{Table}
\else
  \newcommand\tablename{Table}
\fi
}
\@ifpackageloaded{float}{}{\usepackage{float}}
\floatstyle{ruled}
\@ifundefined{c@chapter}{\newfloat{codelisting}{h}{lop}}{\newfloat{codelisting}{h}{lop}[chapter]}
\floatname{codelisting}{Listing}
\newcommand*\listoflistings{\listof{codelisting}{List of Listings}}
\makeatother
\makeatletter
\makeatother
\makeatletter
\@ifpackageloaded{caption}{}{\usepackage{caption}}
\@ifpackageloaded{subcaption}{}{\usepackage{subcaption}}
\makeatother
\usepackage{bookmark}
\IfFileExists{xurl.sty}{\usepackage{xurl}}{} % add URL line breaks if available
\urlstyle{same}
\hypersetup{
  pdftitle={Preparation for Lab 5 from Group 1 (Wörner, Velez, Northe)},
  colorlinks=true,
  linkcolor={blue},
  filecolor={Maroon},
  citecolor={Blue},
  urlcolor={Blue},
  pdfcreator={LaTeX via pandoc}}


\title{Preparation for Lab 5 from Group 1 (Wörner, Velez, Northe)}
\author{}
\date{}
\begin{document}
\maketitle


The source code and the binaries from this report can be found in our
Github
\href{https://github.com/NortheLo/Microwave_Eng_Lab/tree/main/Lab5_Antennas}{repository}.

\subsubsection{Import libs}\label{import-libs}

\begin{Shaded}
\begin{Highlighting}[]
\ImportTok{import}\NormalTok{ numpy }\ImportTok{as}\NormalTok{ np}
\ImportTok{from}\NormalTok{ scipy.constants }\ImportTok{import}\NormalTok{ Boltzmann}
\end{Highlighting}
\end{Shaded}

\subsubsection{Helper constants and
functions}\label{helper-constants-and-functions}

\begin{Shaded}
\begin{Highlighting}[]
\NormalTok{k\_B }\OperatorTok{=}\NormalTok{ Boltzmann     }\CommentTok{\# Boltzman constant}
\NormalTok{c0 }\OperatorTok{=} \FloatTok{3e8} \CommentTok{\#speed of light}

\NormalTok{db2lin }\OperatorTok{=} \KeywordTok{lambda}\NormalTok{ db: }\DecValTok{10}\OperatorTok{**}\NormalTok{(db }\OperatorTok{/} \DecValTok{10}\NormalTok{)}

\NormalTok{lin2db }\OperatorTok{=} \KeywordTok{lambda}\NormalTok{ g: }\DecValTok{10} \OperatorTok{*}\NormalTok{ np.log10(g)}

\NormalTok{lin2dBm }\OperatorTok{=} \KeywordTok{lambda}\NormalTok{ g: lin2db(g}\OperatorTok{/}\FloatTok{1e{-}3}\NormalTok{)}

\KeywordTok{def}\NormalTok{ scientific\_2\_str(value: }\BuiltInTok{float}\NormalTok{, unit: }\BuiltInTok{str} \OperatorTok{=} \StringTok{"Ω"}\NormalTok{) }\OperatorTok{{-}\textgreater{}} \BuiltInTok{str}\NormalTok{:}
\NormalTok{    prefixes }\OperatorTok{=}\NormalTok{ [}
\NormalTok{        (}\FloatTok{1e9}\NormalTok{, }\StringTok{\textquotesingle{}G\textquotesingle{}}\NormalTok{),}
\NormalTok{        (}\FloatTok{1e6}\NormalTok{, }\StringTok{\textquotesingle{}M\textquotesingle{}}\NormalTok{),}
\NormalTok{        (}\FloatTok{1e3}\NormalTok{, }\StringTok{\textquotesingle{}k\textquotesingle{}}\NormalTok{),}
\NormalTok{        (}\DecValTok{1}\NormalTok{, }\StringTok{\textquotesingle{}\textquotesingle{}}\NormalTok{),}
\NormalTok{        (}\FloatTok{1e{-}3}\NormalTok{, }\StringTok{\textquotesingle{}m\textquotesingle{}}\NormalTok{),}
\NormalTok{        (}\FloatTok{1e{-}6}\NormalTok{, }\StringTok{\textquotesingle{}µ\textquotesingle{}}\NormalTok{),}
\NormalTok{        (}\FloatTok{1e{-}9}\NormalTok{, }\StringTok{\textquotesingle{}n\textquotesingle{}}\NormalTok{),}
\NormalTok{        (}\FloatTok{1e{-}12}\NormalTok{, }\StringTok{\textquotesingle{}p\textquotesingle{}}\NormalTok{)    }
\NormalTok{    ]}

    \ControlFlowTok{for}\NormalTok{ factor, prefix }\KeywordTok{in}\NormalTok{ prefixes:}
        \ControlFlowTok{if} \BuiltInTok{abs}\NormalTok{(value) }\OperatorTok{\textgreater{}=}\NormalTok{ factor:}
\NormalTok{            formatted }\OperatorTok{=}\NormalTok{ value }\OperatorTok{/}\NormalTok{ factor}
            \ControlFlowTok{return} \SpecialStringTok{f"}\SpecialCharTok{\{}\NormalTok{formatted}\SpecialCharTok{:.3g\}}\SpecialStringTok{ }\SpecialCharTok{\{}\NormalTok{prefix}\SpecialCharTok{\}\{}\NormalTok{unit}\SpecialCharTok{\}}\SpecialStringTok{"}
    \ControlFlowTok{return} \SpecialStringTok{f"}\SpecialCharTok{\{}\NormalTok{value}\SpecialCharTok{:.3g\}}\SpecialStringTok{ }\SpecialCharTok{\{}\NormalTok{unit}\SpecialCharTok{\}}\SpecialStringTok{"}  \CommentTok{\# fallback for very small values}
\end{Highlighting}
\end{Shaded}

\subsection{5.1 Exercise 1}\label{exercise-1}

Calculate the expected gain of an horn antenna with dimensions \(a^{'}\)
= 2.3 cm, \(b^{'}\) = 12 cm at f = 12 GHz.

From skipt page 10 we have the gain approximation for Horn Antennas with
\$ G = \frac{32a^{'}b^{'}}{\pi \lambda_{0}^{2} }\$ where \(a^{'}\) and
\(b^{'}\) being the dimensions of the extended horn antenna

\begin{Shaded}
\begin{Highlighting}[]
\CommentTok{\#\#\# CONSTANTS }\AlertTok{\#\#\#}
\NormalTok{a }\OperatorTok{=} \FloatTok{2.3e{-}2} \CommentTok{\# meters}
\NormalTok{b }\OperatorTok{=} \FloatTok{12e{-}2} \CommentTok{\# meters}
\NormalTok{f }\OperatorTok{=} \FloatTok{12e9} \CommentTok{\#GHz}
\NormalTok{wavelength }\OperatorTok{=}\NormalTok{ c0 }\OperatorTok{/}\NormalTok{ f}

\CommentTok{\#\# Gain calculation}
\NormalTok{gain }\OperatorTok{=}\NormalTok{ (}\DecValTok{32} \OperatorTok{*}\NormalTok{ a }\OperatorTok{*}\NormalTok{ b ) }\OperatorTok{/}\NormalTok{ (np.pi }\OperatorTok{*}\NormalTok{ np.square(wavelength))}
\NormalTok{gaindb }\OperatorTok{=}\NormalTok{ lin2db(gain)}
\BuiltInTok{print}\NormalTok{(}\SpecialStringTok{f"The approximated gain of the horn antenna with the given dimensions is: }\SpecialCharTok{\{}\NormalTok{gaindb}\SpecialCharTok{:.2f\}}\SpecialStringTok{ dB"}\NormalTok{)}
\end{Highlighting}
\end{Shaded}

\begin{verbatim}
The approximated gain of the horn antenna with the given dimensions is: 16.53 dB
\end{verbatim}

\subsection{5.2 Exercise 2}\label{exercise-2}

You have to design a patch for the frequency f = 2 GHz. What are the
length L and width W of the patch on a Duroid substrate with the
thickness d = 2.5mm and \(\epsilon_{r}\) = 2.2?

The function for calculating the patch lengh is given with:

L = \(\frac{1}{2}\) \(\cdot\) \(\frac{c_{0}}{f \sqrt{\epsilon_{r}}}\)
\(\cdot\) \(\alpha\) = \(\frac{1}{2}\) \(\cdot\)
\(\frac{c_{0}}{f\sqrt{\epsilon_{r}}}\) - 2 \(\Delta L\)

with \(\alpha\) beeing the shortening factor. \(\Delta L\) is
approximated with \(\Delta L\) \(\approx\) \(\frac{d}{2}\).

For the calculation of the patch width, the following function is used:
w = \$ \frac{1}{2} \cdot \frac{c_{0}}{f}
\cdot \sqrt{\frac{2}{\epsilon_{r}+1}}\$

\begin{Shaded}
\begin{Highlighting}[]
\CommentTok{\#\#\# CONSTANTS }\AlertTok{\#\#\#}
\NormalTok{d }\OperatorTok{=} \FloatTok{2.5e{-}3} 
\NormalTok{f }\OperatorTok{=} \FloatTok{2e9} \CommentTok{\#GHz}
\NormalTok{epsilon }\OperatorTok{=} \FloatTok{2.2}
\NormalTok{wavelength }\OperatorTok{=}\NormalTok{ c0 }\OperatorTok{/}\NormalTok{ f}
\NormalTok{deltaL }\OperatorTok{=}\NormalTok{  d }\OperatorTok{/} \DecValTok{2}

\CommentTok{\#\# Patch length calculation}
\NormalTok{length\_m }\OperatorTok{=}\NormalTok{ ((}\FloatTok{0.5} \OperatorTok{*}\NormalTok{ c0) }\OperatorTok{/}\NormalTok{ (f }\OperatorTok{*}\NormalTok{ np.sqrt(epsilon))) }\OperatorTok{{-}} \DecValTok{2} \OperatorTok{*}\NormalTok{ deltaL}
\NormalTok{length\_mm }\OperatorTok{=}\NormalTok{ length\_m }\OperatorTok{*} \DecValTok{100}

\CommentTok{\#\# Patch width calculation}
\NormalTok{width }\OperatorTok{=}\NormalTok{ ((}\FloatTok{0.5} \OperatorTok{*}\NormalTok{ c0) }\OperatorTok{/}\NormalTok{ f ) }\OperatorTok{*}\NormalTok{ np.sqrt(}\DecValTok{2} \OperatorTok{/}\NormalTok{ (epsilon}\OperatorTok{+}\DecValTok{1}\NormalTok{))}
\NormalTok{width\_mm }\OperatorTok{=}\NormalTok{ width }\OperatorTok{*} \DecValTok{100}

\BuiltInTok{print}\NormalTok{(}\SpecialStringTok{f"The approximated length of the antenna patch with the given conditions is: }\SpecialCharTok{\{}\NormalTok{length\_mm}\SpecialCharTok{:.3f\}}\SpecialStringTok{ mm"}\NormalTok{)}
\BuiltInTok{print}\NormalTok{(}\SpecialStringTok{f"The approximated width of the antenna patch with the given conditions is: }\SpecialCharTok{\{}\NormalTok{width\_mm}\SpecialCharTok{:.3f\}}\SpecialStringTok{ mm"}\NormalTok{)}
\end{Highlighting}
\end{Shaded}

\begin{verbatim}
The approximated length of the antenna patch with the given conditions is: 4.806 mm
The approximated width of the antenna patch with the given conditions is: 5.929 mm
\end{verbatim}

\subsection{5.3 Exercise 3}\label{exercise-3}

You have to design a phase shifter for a phase-driven array. Three
pivoting angles are needed: \(\theta\) = 90°, \(\theta\) = 70°,
\(\theta\) = 50°. - Calculate the needed phase increments. - Calculate
the length of the needed coaxial lines (Teflon fill, \(\epsilon_{r}\) =
2.2), to use them as phase shifters. The distance between the elements
is d = 75mm and the frequency is f = 2.0 GHz.

To calculate the needed phase increments, we use the formular \$
\Delta \phi = - \frac{2\pi}{\lambda} \cdot d \cdot cos \theta \$, where
d is the distance between two antennas.

\begin{Shaded}
\begin{Highlighting}[]
\CommentTok{\#\#\# CONSTANTS }\AlertTok{\#\#\#}
\NormalTok{d }\OperatorTok{=} \FloatTok{75e{-}3} \CommentTok{\#\#75mm}
\NormalTok{theta1 }\OperatorTok{=} \DecValTok{90} \CommentTok{\#90 degree}
\NormalTok{theta2 }\OperatorTok{=} \DecValTok{70} \CommentTok{\#70 degree}
\NormalTok{theta3 }\OperatorTok{=} \DecValTok{50} \CommentTok{\#50 degree}
\NormalTok{epsilon }\OperatorTok{=} \FloatTok{2.2}
\NormalTok{f }\OperatorTok{=} \FloatTok{2e9} \CommentTok{\#2 GHz}
\NormalTok{wavelength }\OperatorTok{=}\NormalTok{ c0 }\OperatorTok{/}\NormalTok{ f}


\CommentTok{\#\#\# Calculation of the needed phase increment}

\CommentTok{\#\# radiation direction theta1: }
\NormalTok{phase\_inc\_1 }\OperatorTok{=}\NormalTok{ ((}\OperatorTok{{-}}\DecValTok{2}\OperatorTok{*}\NormalTok{ np.pi) }\OperatorTok{/}\NormalTok{ wavelength) }\OperatorTok{*}\NormalTok{ d }\OperatorTok{*}\NormalTok{ np.cos(np.deg2rad(theta1))}

\CommentTok{\#\# radiation direction theta2: }
\NormalTok{phase\_inc\_2 }\OperatorTok{=}\NormalTok{ ((}\OperatorTok{{-}}\DecValTok{2}\OperatorTok{*}\NormalTok{ np.pi) }\OperatorTok{/}\NormalTok{ wavelength) }\OperatorTok{*}\NormalTok{ d }\OperatorTok{*}\NormalTok{ np.cos(np.deg2rad(theta2))}

\CommentTok{\#\# radiation direction theta3: }
\NormalTok{phase\_inc\_3 }\OperatorTok{=}\NormalTok{ ((}\OperatorTok{{-}}\DecValTok{2}\OperatorTok{*}\NormalTok{ np.pi) }\OperatorTok{/}\NormalTok{ wavelength) }\OperatorTok{*}\NormalTok{ d }\OperatorTok{*}\NormalTok{ np.cos(np.deg2rad(theta3))}

\BuiltInTok{print}\NormalTok{(}\SpecialStringTok{f"The needed phase increment for angle theta 1 is: }\SpecialCharTok{\{}\NormalTok{phase\_inc\_1}\SpecialCharTok{:.3f\}}\SpecialStringTok{."}\NormalTok{)}
\BuiltInTok{print}\NormalTok{(}\SpecialStringTok{f"The needed phase increment for angle theta 2 is: }\SpecialCharTok{\{}\NormalTok{phase\_inc\_2}\SpecialCharTok{:.3f\}}\SpecialStringTok{."}\NormalTok{)}
\BuiltInTok{print}\NormalTok{(}\SpecialStringTok{f"The needed phase increment for angle theta 3 is: }\SpecialCharTok{\{}\NormalTok{phase\_inc\_3}\SpecialCharTok{:.3f\}}\SpecialStringTok{."}\NormalTok{)}
\end{Highlighting}
\end{Shaded}

\begin{verbatim}
The needed phase increment for angle theta 1 is: -0.000.
The needed phase increment for angle theta 2 is: -1.074.
The needed phase increment for angle theta 3 is: -2.019.
\end{verbatim}

In the following, we calculate the length of the needed coaxial lines,
to use them as phase shifters. We calculate the length with the
resulting phase increments from the task before and the equation \$ l =
\frac{\Delta \phi}{2 \pi}
\cdot \frac{c_{0}}{f \cdot \sqrt{\epsilon_{r}} } \$. Since a negative
phase increment from left to right is a positive increment from right to
left, we can just use the absolute value of the calculated length.

\begin{Shaded}
\begin{Highlighting}[]
\CommentTok{\#\#\# Calculation of the needed line length to shift the phase by 0 to get the pivoting angle theta1: }
\NormalTok{length1 }\OperatorTok{=} \BuiltInTok{abs}\NormalTok{((phase\_inc\_1 }\OperatorTok{/}\NormalTok{ (}\DecValTok{2} \OperatorTok{*}\NormalTok{ np.pi)) }\OperatorTok{*}\NormalTok{ (wavelength }\OperatorTok{/}\NormalTok{ np.sqrt(epsilon))) }\OperatorTok{*} \DecValTok{100}

\CommentTok{\#\#\# Calculation of the needed line length to shift the phase by 1 to get the pivoting angle theta2: }
\NormalTok{length2 }\OperatorTok{=} \BuiltInTok{abs}\NormalTok{((phase\_inc\_2 }\OperatorTok{/}\NormalTok{ (}\DecValTok{2} \OperatorTok{*}\NormalTok{ np.pi)) }\OperatorTok{*}\NormalTok{ (wavelength }\OperatorTok{/}\NormalTok{ np.sqrt(epsilon))) }\OperatorTok{*} \DecValTok{100}

\CommentTok{\#\#\# Calculation of the needed line length to shift the phase by 2 to get the pivoting angle theta3: }
\NormalTok{length3 }\OperatorTok{=} \BuiltInTok{abs}\NormalTok{((phase\_inc\_3 }\OperatorTok{/}\NormalTok{ (}\DecValTok{2} \OperatorTok{*}\NormalTok{ np.pi)) }\OperatorTok{*}\NormalTok{ (wavelength }\OperatorTok{/}\NormalTok{ np.sqrt(epsilon))) }\OperatorTok{*} \DecValTok{100}

\BuiltInTok{print}\NormalTok{(}\SpecialStringTok{f"The necessary length of the wire to shift the phase far enough to reach pivoting angle theta1 is: }\SpecialCharTok{\{}\NormalTok{length1}\SpecialCharTok{:.3f\}}\SpecialStringTok{ mm."}\NormalTok{)}
\BuiltInTok{print}\NormalTok{(}\SpecialStringTok{f"The necessary length of the wire to shift the phase far enough to reach pivoting angle theta2 is: }\SpecialCharTok{\{}\NormalTok{length2}\SpecialCharTok{:.3f\}}\SpecialStringTok{ mm."}\NormalTok{)}
\BuiltInTok{print}\NormalTok{(}\SpecialStringTok{f"The necessary length of the wire to shift the phase far enough to reach pivoting angle theta3 is: }\SpecialCharTok{\{}\NormalTok{length3}\SpecialCharTok{:.3f\}}\SpecialStringTok{ mm."}\NormalTok{)}
\end{Highlighting}
\end{Shaded}

\begin{verbatim}
The necessary length of the wire to shift the phase far enough to reach pivoting angle theta1 is: 0.000 mm.
The necessary length of the wire to shift the phase far enough to reach pivoting angle theta2 is: 1.729 mm.
The necessary length of the wire to shift the phase far enough to reach pivoting angle theta3 is: 3.250 mm.
\end{verbatim}

\subsection{5.4 Exercise 4}\label{exercise-4}

The power radiation pattern for a horn antenna extended in the E-plane
has to be calculated in the E-plane ( \(\theta\) = 0) for angles from
\(\theta\) = 0° to \(\theta\) = 180° and plotted in polar cooradinates.
- State the formulas to calculate the amplitudes in the far field and
plot them in 1°-steps in a Matlab plot (in dB). - State the needed
formulas and calculate the following values. Also, draw them by hand in
your Matlab plot. -- The first 4 nulls -- The half power beam width
(approximative formula) Dimensions of the antenna aperture: a = 2:3 cm,
b = 12 cm, frequency f = 11.85 GHz.

\subsubsection{E-plane Radiation Pattern with Nulls and
HPBW}\label{e-plane-radiation-pattern-with-nulls-and-hpbw}

This nex plot shows the normalized far-field radiation pattern of a horn
antenna extended in the E-plane.\\
The \textbf{red circles} mark the first four nulls w.r.t. the main lobe
at 90°.\\
The \textbf{green dots} indicate the \textbf{Half Power Beam Width
(HPBW)} --- the angular region where the radiated power is within 3 dB
of the maximum (half power). This region represents the main lobe of the
antenna's radiation pattern.

\begin{Shaded}
\begin{Highlighting}[]
\ImportTok{import}\NormalTok{ numpy }\ImportTok{as}\NormalTok{ np}
\ImportTok{import}\NormalTok{ matplotlib.pyplot }\ImportTok{as}\NormalTok{ plt}

\CommentTok{\# {-}{-}{-} Constants {-}{-}{-}}
\NormalTok{c }\OperatorTok{=} \FloatTok{3e8}  \CommentTok{\# speed of light in m/s}
\NormalTok{f }\OperatorTok{=} \FloatTok{11.85e9}  \CommentTok{\# frequency in Hz}
\NormalTok{λ }\OperatorTok{=}\NormalTok{ c }\OperatorTok{/}\NormalTok{ f  }\CommentTok{\# wavelength in meters}
\NormalTok{b\_prime }\OperatorTok{=} \FloatTok{0.12}  \CommentTok{\# aperture size in the E{-}plane (in meters)}
\NormalTok{beta }\OperatorTok{=} \DecValTok{2} \OperatorTok{*}\NormalTok{ np.pi }\OperatorTok{/}\NormalTok{ λ  }\CommentTok{\# wave number}

\CommentTok{\# {-}{-}{-} Angular range (0° to 180°) {-}{-}{-}}
\NormalTok{theta\_deg }\OperatorTok{=}\NormalTok{ np.arange(}\DecValTok{0}\NormalTok{, }\DecValTok{181}\NormalTok{, }\DecValTok{1}\NormalTok{)}
\NormalTok{theta\_rad }\OperatorTok{=}\NormalTok{ np.radians(theta\_deg)}

\CommentTok{\# {-}{-}{-} Compute B = (b\textquotesingle{} * beta * cos(theta)) / 2 {-}{-}{-}}
\NormalTok{B }\OperatorTok{=}\NormalTok{ (b\_prime }\OperatorTok{*}\NormalTok{ beta }\OperatorTok{*}\NormalTok{ np.cos(theta\_rad)) }\OperatorTok{/} \DecValTok{2}
\NormalTok{B[B }\OperatorTok{==} \DecValTok{0}\NormalTok{] }\OperatorTok{=} \FloatTok{1e{-}12}  \CommentTok{\# Avoid division by zero}

\CommentTok{\# {-}{-}{-} Compute normalized far{-}field amplitude |U(theta)/Umax| = |sin(B)/B| {-}{-}{-}}
\NormalTok{U\_norm }\OperatorTok{=}\NormalTok{ np.}\BuiltInTok{abs}\NormalTok{(np.sin(B) }\OperatorTok{/}\NormalTok{ B)}
\NormalTok{U\_dB }\OperatorTok{=} \DecValTok{20} \OperatorTok{*}\NormalTok{ np.log10(U\_norm }\OperatorTok{/}\NormalTok{ np.}\BuiltInTok{max}\NormalTok{(U\_norm))  }\CommentTok{\# normalized in dB}

\CommentTok{\# {-}{-}{-} Compute approximate HPBW {-}{-}{-}}
\NormalTok{HPBW }\OperatorTok{=} \DecValTok{51} \OperatorTok{*}\NormalTok{ λ }\OperatorTok{/}\NormalTok{ b\_prime  }\CommentTok{\# in degrees}

\CommentTok{\# {-}{-}{-} Find first 4 nulls: B = nπ {-}{-}\textgreater{} cos(θ\_null) {-}{-}{-}}
\NormalTok{null\_angles }\OperatorTok{=}\NormalTok{ []}
\ControlFlowTok{for}\NormalTok{ n }\KeywordTok{in} \BuiltInTok{range}\NormalTok{(}\DecValTok{1}\NormalTok{, }\DecValTok{5}\NormalTok{):}
\NormalTok{    B\_null }\OperatorTok{=}\NormalTok{ n }\OperatorTok{*}\NormalTok{ np.pi}
\NormalTok{    cos\_theta\_null }\OperatorTok{=} \DecValTok{2} \OperatorTok{*}\NormalTok{ B\_null }\OperatorTok{/}\NormalTok{ (b\_prime }\OperatorTok{*}\NormalTok{ beta)}
    \ControlFlowTok{if}\NormalTok{ np.}\BuiltInTok{abs}\NormalTok{(cos\_theta\_null) }\OperatorTok{\textless{}=} \DecValTok{1}\NormalTok{:}
\NormalTok{        theta\_null }\OperatorTok{=}\NormalTok{ np.degrees(np.arccos(cos\_theta\_null))}
\NormalTok{        null\_angles.append(theta\_null)}

\CommentTok{\# {-}{-}{-} Plot in polar coordinates {-}{-}{-}}
\NormalTok{fig }\OperatorTok{=}\NormalTok{ plt.figure(figsize}\OperatorTok{=}\NormalTok{(}\DecValTok{8}\NormalTok{, }\DecValTok{6}\NormalTok{))}
\NormalTok{ax }\OperatorTok{=}\NormalTok{ fig.add\_subplot(}\DecValTok{111}\NormalTok{, polar}\OperatorTok{=}\VariableTok{True}\NormalTok{)}

\CommentTok{\# Plot radiation pattern}
\NormalTok{ax.plot(np.radians(theta\_deg), U\_dB, label}\OperatorTok{=}\StringTok{"Normalized |E(θ)/Emax| (dB)"}\NormalTok{)}

\CommentTok{\# Mark nulls}
\ControlFlowTok{for}\NormalTok{ angle }\KeywordTok{in}\NormalTok{ null\_angles:}
\NormalTok{    ax.plot(np.radians(angle), [}\OperatorTok{{-}}\DecValTok{40}\NormalTok{], }\StringTok{\textquotesingle{}ro\textquotesingle{}}\NormalTok{)}
\NormalTok{    ax.text(np.radians(angle), }\OperatorTok{{-}}\DecValTok{38}\NormalTok{, }\SpecialStringTok{f\textquotesingle{}}\SpecialCharTok{\{}\NormalTok{angle}\SpecialCharTok{:.1f\}}\SpecialStringTok{°\textquotesingle{}}\NormalTok{, color}\OperatorTok{=}\StringTok{\textquotesingle{}red\textquotesingle{}}\NormalTok{, ha}\OperatorTok{=}\StringTok{\textquotesingle{}center\textquotesingle{}}\NormalTok{, fontsize}\OperatorTok{=}\DecValTok{8}\NormalTok{)}

\CommentTok{\# Mark {-}3 dB points (HPBW limits)}
\NormalTok{half\_beam }\OperatorTok{=}\NormalTok{ HPBW }\OperatorTok{/} \DecValTok{2}
\NormalTok{left\_bound }\OperatorTok{=} \DecValTok{90} \OperatorTok{{-}}\NormalTok{ half\_beam}
\NormalTok{right\_bound }\OperatorTok{=} \DecValTok{90} \OperatorTok{+}\NormalTok{ half\_beam}
\NormalTok{ax.plot(np.radians(left\_bound), [}\OperatorTok{{-}}\DecValTok{3}\NormalTok{], }\StringTok{\textquotesingle{}go\textquotesingle{}}\NormalTok{, label}\OperatorTok{=}\StringTok{\textquotesingle{}{-}3 dB points\textquotesingle{}}\NormalTok{)}
\NormalTok{ax.plot(np.radians(right\_bound), [}\OperatorTok{{-}}\DecValTok{3}\NormalTok{], }\StringTok{\textquotesingle{}go\textquotesingle{}}\NormalTok{)}

\CommentTok{\# Shade HPBW region in light green}
\NormalTok{theta\_fill }\OperatorTok{=}\NormalTok{ np.radians(np.arange(}\BuiltInTok{int}\NormalTok{(left\_bound), }\BuiltInTok{int}\NormalTok{(right\_bound) }\OperatorTok{+} \DecValTok{1}\NormalTok{))}
\NormalTok{radius\_fill }\OperatorTok{=}\NormalTok{ np.interp(theta\_fill, np.radians(theta\_deg), U\_dB)}
\NormalTok{ax.fill\_between(theta\_fill, radius\_fill, }\OperatorTok{{-}}\DecValTok{40}\NormalTok{, color}\OperatorTok{=}\StringTok{\textquotesingle{}lightgreen\textquotesingle{}}\NormalTok{, alpha}\OperatorTok{=}\FloatTok{0.5}\NormalTok{, label}\OperatorTok{=}\StringTok{\textquotesingle{}HPBW region\textquotesingle{}}\NormalTok{)}

\CommentTok{\# Final plot settings}
\NormalTok{ax.set\_title(}\StringTok{"E{-}plane Radiation Pattern of Horn Antenna"}\NormalTok{, va}\OperatorTok{=}\StringTok{\textquotesingle{}bottom\textquotesingle{}}\NormalTok{)}
\NormalTok{ax.set\_rticks([}\OperatorTok{{-}}\DecValTok{40}\NormalTok{, }\OperatorTok{{-}}\DecValTok{30}\NormalTok{, }\OperatorTok{{-}}\DecValTok{20}\NormalTok{, }\OperatorTok{{-}}\DecValTok{10}\NormalTok{, }\DecValTok{0}\NormalTok{])}
\NormalTok{ax.set\_rlabel\_position(}\DecValTok{180}\NormalTok{)}
\NormalTok{ax.grid(}\VariableTok{True}\NormalTok{)}
\NormalTok{ax.legend(loc}\OperatorTok{=}\StringTok{\textquotesingle{}lower right\textquotesingle{}}\NormalTok{)}
\NormalTok{plt.show()}

\CommentTok{\# {-}{-}{-} Print Results {-}{-}{-}}
\BuiltInTok{print}\NormalTok{(}\StringTok{"Wavelength λ = }\SpecialCharTok{\{:.3f\}}\StringTok{ cm"}\NormalTok{.}\BuiltInTok{format}\NormalTok{(λ }\OperatorTok{*} \DecValTok{100}\NormalTok{))}
\BuiltInTok{print}\NormalTok{(}\StringTok{"Approximate Half Power Beam Width (HPBW) ≈ }\SpecialCharTok{\{:.2f\}}\StringTok{°"}\NormalTok{.}\BuiltInTok{format}\NormalTok{(HPBW))}
\BuiltInTok{print}\NormalTok{(}\StringTok{"First 4 nulls at angles (deg):"}\NormalTok{, [}\SpecialStringTok{f"}\SpecialCharTok{\{}\NormalTok{a}\SpecialCharTok{:.2f\}}\SpecialStringTok{"} \ControlFlowTok{for}\NormalTok{ a }\KeywordTok{in}\NormalTok{ null\_angles])}
\end{Highlighting}
\end{Shaded}

\pandocbounded{\includegraphics[keepaspectratio]{Lab5_Antennas_Preliminary_Group1_files/figure-pdf/cell-8-output-1.png}}

\begin{verbatim}
Wavelength λ = 2.532 cm
Approximate Half Power Beam Width (HPBW) ≈ 10.76°
First 4 nulls at angles (deg): ['77.82', '65.04', '50.73', '32.45']
\end{verbatim}

\subsection{5.5 Exercise 5}\label{exercise-5}

You are building a 13 GHz communication link between two office
buildings with two identical antennas. The main beams of both antennas
are oriented against each other. Both antennas have a gain of 9 dBi. The
transmit power \(PT_{x}\) is 6dBm. The length of the communication link
is 100m. Calculate the receive power \(PR_{x}\). You observe that the
system is not working. You hear from the manufacturer of the receiver
that the receive power must be at least \(PR_{x}\) = -60dBm for the
communication link to work. To reach this receive power, only one
antenna shall be exchanged. Calculate the gain of the new antenna (in
dBi).

\subsubsection{Part 1: Calculate Received
Power}\label{part-1-calculate-received-power}

This section uses the Friis transmission equation to calculate the
received power at 13\,GHz over a 100\,m link using two antennas with
9\,dBi gain. From the result, the minimum receive power is not achieved:

\begin{Shaded}
\begin{Highlighting}[]
\CommentTok{\# {-}{-}{-} Part 1: Compute received power with both antennas having 9 dBi gain {-}{-}{-}}

\CommentTok{\# Given values}
\NormalTok{c }\OperatorTok{=} \FloatTok{3e8}  \CommentTok{\# speed of light [m/s]}
\NormalTok{f }\OperatorTok{=} \FloatTok{13e9}  \CommentTok{\# frequency [Hz]}
\NormalTok{λ }\OperatorTok{=}\NormalTok{ c }\OperatorTok{/}\NormalTok{ f  }\CommentTok{\# wavelength [m]}
\NormalTok{r }\OperatorTok{=} \DecValTok{100}  \CommentTok{\# distance [m]}
\NormalTok{Pt\_dBm }\OperatorTok{=} \DecValTok{6}  \CommentTok{\# transmit power [dBm]}
\NormalTok{G1\_dBi }\OperatorTok{=} \DecValTok{9}  \CommentTok{\# gain of antenna 1 [dBi]}
\NormalTok{G2\_dBi }\OperatorTok{=} \DecValTok{9}  \CommentTok{\# gain of antenna 2 [dBi]}

\CommentTok{\# Friis transmission formula in dB}
\NormalTok{path\_loss\_dB }\OperatorTok{=} \DecValTok{20} \OperatorTok{*}\NormalTok{ np.log10(λ }\OperatorTok{/}\NormalTok{ (}\DecValTok{4} \OperatorTok{*}\NormalTok{ np.pi }\OperatorTok{*}\NormalTok{ r))}
\NormalTok{Pr\_dBm }\OperatorTok{=}\NormalTok{ Pt\_dBm }\OperatorTok{+}\NormalTok{ G1\_dBi }\OperatorTok{+}\NormalTok{ G2\_dBi }\OperatorTok{+}\NormalTok{ path\_loss\_dB}

\CommentTok{\# Output}
\BuiltInTok{print}\NormalTok{(}\StringTok{"=== PART 1 ==="}\NormalTok{)}
\BuiltInTok{print}\NormalTok{(}\SpecialStringTok{f"Wavelength λ = }\SpecialCharTok{\{λ}\OperatorTok{*}\DecValTok{100}\SpecialCharTok{:.2f\}}\SpecialStringTok{ cm"}\NormalTok{)}
\BuiltInTok{print}\NormalTok{(}\SpecialStringTok{f"Path Loss = }\SpecialCharTok{\{}\NormalTok{path\_loss\_dB}\SpecialCharTok{:.2f\}}\SpecialStringTok{ dB"}\NormalTok{)}
\BuiltInTok{print}\NormalTok{(}\SpecialStringTok{f"Received Power = }\SpecialCharTok{\{}\NormalTok{Pr\_dBm}\SpecialCharTok{:.2f\}}\SpecialStringTok{ dBm"}\NormalTok{)}
\end{Highlighting}
\end{Shaded}

\begin{verbatim}
=== PART 1 ===
Wavelength λ = 2.31 cm
Path Loss = -94.72 dB
Received Power = -70.72 dBm
\end{verbatim}

\subsubsection{Part 2: Determine Required
Gain}\label{part-2-determine-required-gain}

This next section calculates the minimum gain needed for a replacement
antenna to meet the link budget.

\begin{Shaded}
\begin{Highlighting}[]
\CommentTok{\# {-}{-}{-} Part 2: Compute required gain of new antenna if system doesn\textquotesingle{}t meet threshold {-}{-}{-}}

\NormalTok{Rx\_min\_dBm }\OperatorTok{=} \OperatorTok{{-}}\DecValTok{60}  \CommentTok{\# minimum receive power for link to work}

\CommentTok{\# Compute required gain of one antenna if only one is replaced}
\NormalTok{G\_new\_dBi }\OperatorTok{=}\NormalTok{ Rx\_min\_dBm }\OperatorTok{{-}}\NormalTok{ Pr\_dBm }\OperatorTok{+}\NormalTok{ G1\_dBi}
\BuiltInTok{print}\NormalTok{(}\StringTok{"}\CharTok{\textbackslash{}n}\StringTok{=== PART 2 ==="}\NormalTok{)}
\BuiltInTok{print}\NormalTok{(}\SpecialStringTok{f"Minimum required gain of replacement antenna: }\SpecialCharTok{\{}\NormalTok{G\_new\_dBi}\SpecialCharTok{:.2f\}}\SpecialStringTok{ dBi"}\NormalTok{)}

\CommentTok{\# Optional: verify by recalculating received power}
\NormalTok{Pr\_dBm\_check }\OperatorTok{=}\NormalTok{ Pt\_dBm }\OperatorTok{+}\NormalTok{ G1\_dBi }\OperatorTok{+}\NormalTok{ G\_new\_dBi }\OperatorTok{+}\NormalTok{ path\_loss\_dB}
\BuiltInTok{print}\NormalTok{(}\SpecialStringTok{f"Recomputed received power with new antenna: }\SpecialCharTok{\{}\NormalTok{Pr\_dBm\_check}\SpecialCharTok{:.2f\}}\SpecialStringTok{ dBm"}\NormalTok{)}
\end{Highlighting}
\end{Shaded}

\begin{verbatim}

=== PART 2 ===
Minimum required gain of replacement antenna: 19.72 dBi
Recomputed received power with new antenna: -60.00 dBm
\end{verbatim}




\end{document}
